\subsection{Discrete random variables}

\subsubsection{Definition of random variables}

A random variable is a function from the sample space to the real numbers.  
It assigns to each outcome of the experiment a numerical value, determined once the outcome is known.

\bigskip\textsc{Example: Weight of a student}

Let the sample space be \(\Omega = \{a,b,c,d\}\), where each element is a student.  
The experiment is to pick one student at random and record their weight in kilograms.  
Define \(W(\omega)\) as the weight of student \(\omega\).  
Then \(W\) is a random variable, because for each outcome \(\omega \in \Omega\) it produces a real number \(W(\omega)\).

\bigskip\textsc{Another random variable and a function of them}

On the same experiment, define \(H(\omega)\) as the height (in meters) of student \(\omega\).  
This \(H\) is also a random variable.  
We can form a new random variable by combining them, for instance the body mass index
\[
B(\omega) = \frac{W(\omega)}{H(\omega)^2}.
\]
Once the outcome \(\omega\) is known, both \(W(\omega)\) and \(H(\omega)\) are known, so \(B(\omega)\) is determined as well.

\bigskip\textsc{Discrete and continuous random variables}

Random variables can be discrete or continuous.  
For example, the number of heads in \(10\) coin tosses is a discrete random variable taking values in \(\{0,1,\dots,10\}\).  
The time at which an event occurs, measured with infinite accuracy, is modeled as a continuous random variable taking real values in some interval.

\bigskip\textsc{Notation}

We use uppercase letters such as \(X, Y, W, H\) for random variables (the functions) and lowercase letters such as \(x, y, w, h\) for their numerical values.  
If \(X\) and \(Y\) are random variables, then \(X+Y\) is also a random variable, defined by
\[
(X+Y)(\omega) = X(\omega) + Y(\omega).
\]
In general, any function of random variables is again a random variable.



\subsubsection{Probability mass function}

A discrete random variable can take different numerical values, and some of these values are more likely than others.  
These relative likelihoods are described by the \textit{probability mass function} (PMF), which gives the probability of each possible numerical value.

\bigskip\textsc{Definition of the PMF}

Let \(X\) be a discrete random variable.  
For any real number \(x\), we write
\[
p_X(x) = \mathbb{P}(X = x).
\]
This is the probability of the event that \(X\) takes the value \(x\).  
The function \(p_X\) is called the PMF (or probability law, or probability distribution) of \(X\).  

The argument \(x\) ranges over the possible values that \(X\) can take.  
For each such \(x\), \(p_X(x)\) is a number between \(0\) and \(1\).  
If we had another random variable \(Y\) on the same sample space, its PMF would be denoted \(p_Y(y) = \mathbb{P}(Y = y)\).

\bigskip\textsc{Basic properties}

The PMF satisfies two fundamental properties:

\[
p_X(x) \ge 0 \quad \text{for all } x,
\]
\[
\sum_{x} p_X(x) = 1,
\]
where the sum is over all possible values of \(X\).  
The first property holds because probabilities are never negative.  
The second holds because the events \(\{X = x\}\) for all possible \(x\) are disjoint and together cover the entire sample space, so their probabilities must add up to \(1\).

\bigskip\textsc{Example: A simple random variable}

Consider an experiment with four equally likely outcomes \(a, b, c, d\).  
Suppose a random variable \(X\) is defined by
\[
X(a) = 5, \quad X(b) = 5, \quad X(c) = 4, \quad X(d) = 3.
\]
Each outcome has probability \(1/4\).  

The event \(\{X = 5\}\) is \(\{a, b\}\), so
\[
p_X(5) = \mathbb{P}(X = 5) = \mathbb{P}(\{a, b\}) = \frac{1}{4} + \frac{1}{4} = \frac{1}{2}.
\]
Similarly,
\[
p_X(4) = \mathbb{P}(X = 4) = \mathbb{P}(\{c\}) = \frac{1}{4},
\]
\[
p_X(3) = \mathbb{P}(X = 3) = \mathbb{P}(\{d\}) = \frac{1}{4}.
\]
For any other \(x\), we have \(p_X(x) = 0\).  
These values satisfy \(p_X(3) + p_X(4) + p_X(5) = \frac{1}{4} + \frac{1}{4} + \frac{1}{2} = 1\), as required.

\bigskip\textsc{Example: Sum of two tetrahedral dice}

Consider rolling two independent tetrahedral dice, each with faces \(\{1,2,3,4\}\).  
Let \(X\) be the result of the first roll and \(Y\) the result of the second roll.  
All \(16\) pairs \((i,j)\) with \(i,j \in \{1,2,3,4\}\) are assumed equally likely, each with probability \(1/16\).  

Define a new random variable
\[
Z = X + Y.
\]
The possible values of \(Z\) are \(2,3,4,5,6,7,8\).  
To compute the PMF of \(Z\), we proceed value by value, each time collecting the outcomes that give that sum and adding their probabilities.

For example, the event \(\{Z = 2\}\) is \(\{(1,1)\}\), so
\[
p_Z(2) = \mathbb{P}(Z = 2) = \frac{1}{16}.
\]
The event \(\{Z = 3\}\) is \(\{(1,2), (2,1)\}\), so
\[
p_Z(3) = \mathbb{P}(Z = 3) = \frac{2}{16}.
\]
The event \(\{Z = 4\}\) is \(\{(1,3), (2,2), (3,1)\}\), so
\[
p_Z(4) = \mathbb{P}(Z = 4) = \frac{3}{16}.
\]
Continuing in the same way, we find
\[
p_Z(5) = \frac{4}{16}, \quad
p_Z(6) = \frac{3}{16}, \quad
p_Z(7) = \frac{2}{16}, \quad
p_Z(8) = \frac{1}{16}.
\]

These numbers are all nonnegative and they add up to \(1\):
\[
\frac{1}{16} + \frac{2}{16} + \frac{3}{16} + \frac{4}{16} + \frac{3}{16} + \frac{2}{16} + \frac{1}{16} = 1.
\]
Thus we have completely determined the probability mass function of \(Z\).



\subsubsection{Bernoulli and indicator random variables}

A \textit{Bernoulli} random variable is a discrete random variable that takes only the values \(0\) and \(1\).  
Its distribution is determined by a parameter \(p \in [0,1]\), interpreted as the probability of the outcome \(1\).  

Using PMF notation, a Bernoulli random variable \(X\) satisfies
\[
\mathbb{P}(X = 1) = p, \qquad \mathbb{P}(X = 0) = 1 - p.
\]
If we plot its PMF, we obtain two bars: one at \(x=0\) of height \(1-p\), and one at \(x=1\) of height \(p\).  
Bernoulli random variables model experiments with two possible outcomes, such as success/failure or heads/tails.

\bigskip\textsc{Indicator random variables}

Let \(\Omega\) be a sample space and let \(A \subseteq \Omega\) be an event.  
The \textit{indicator} random variable of \(A\) is defined by
\[
I_A(\omega) =
\begin{cases}
1 & \text{if } \omega \in A,\\
0 & \text{if } \omega \notin A.
\end{cases}
\]
Thus \(I_A = 1\) exactly when the event \(A\) occurs, and \(I_A = 0\) otherwise.  

Since \(I_A\) takes only the values \(0\) and \(1\), it is a Bernoulli random variable.  
Moreover,
\[
\mathbb{P}(I_A = 1) = \mathbb{P}(A), \qquad \mathbb{P}(I_A = 0) = 1 - \mathbb{P}(A),
\]
so \(I_A\) is Bernoulli with parameter \(p = \mathbb{P}(A)\).  
Indicator random variables are useful because they convert questions about events into questions about random variables, where algebraic manipulations are often easier.


\subsubsection{Uniform random variables}

A \textit{discrete uniform} random variable is one whose possible values all have the same probability.  
It takes integer values in some finite interval and assigns equal mass to each of them.

\bigskip\textsc{Definition}

Let \(a\) and \(b\) be integers with \(a \le b\).  
We consider the experiment of choosing an integer at random from the set
\[
\{a, a+1, a+2, \dots, b\}.
\]
There are \(b - a + 1\) possible values.  

Let \(X\) be the random variable that equals the chosen integer.  
Then \(X\) has PMF
\[
\mathbb{P}(X = x) =
\begin{cases}
\dfrac{1}{b - a + 1} & \text{if } x \in \{a, a+1, \dots, b\},\\[4pt]
0 & \text{otherwise.}
\end{cases}
\]
In this setting, the outcome of the experiment is already a number, and this number is the value of the random variable, so there is no distinction between outcome and numerical value.

\bigskip\textsc{Example: Seconds on a digital clock}

Consider the seconds reading on a digital clock.  
The possible values are the integers from \(0\) to \(59\), so there are \(60\) possibilities.  
If you glance at the clock at a time chosen completely at random, there is no reason to think one second value is more likely than another.  

Let \(S\) be the seconds reading.  
Then \(S\) is (ideally) discrete uniform on \(\{0,1,\dots,59\}\), with
\[
\mathbb{P}(S = s) = \frac{1}{60}, \quad s = 0,1,\dots,59.
\]

\bigskip\textsc{Degenerate case}

If \(a = b\), then the range of \(X\) consists of a single value \(a\).  
In that case,
\[
\mathbb{P}(X = a) = 1,
\]
and \(X\) is a \textit{constant} or \textit{degenerate} random variable.  
It has no randomness in the usual sense, but it still fits the formal definition of a random variable as a function on the sample space.


\subsubsection{Binomial random variables}

A \textit{binomial} random variable models the number of successes in a fixed number of independent, identical trials, where each trial has two outcomes (success/failure) and the same success probability \(p\).  
Typical examples are repeated coin tosses or repeated yes/no experiments.

\bigskip\textsc{Definition}

Consider \(n\) independent trials.  
On each trial, the probability of success is \(p\) and the probability of failure is \(1-p\).  
Define the random variable \(X\) to be the number of successes observed in these \(n\) trials.  

Then \(X\) is called a \textit{binomial} random variable with parameters \(n\) and \(p\).  
Its possible values are \(0,1,\dots,n\).

\bigskip\textsc{PMF of a binomial random variable}

We now recall the probability that \(X\) takes a specific value \(k\).  
The event \(\{X = k\}\) consists of all sequences of \(n\) trials that contain exactly \(k\) successes and \(n-k\) failures.  
Each particular sequence with \(k\) successes has probability
\[
p^k (1-p)^{n-k},
\]
because we multiply a factor \(p\) for each success and a factor \(1-p\) for each failure.  

The number of such sequences is
\[
\binom{n}{k},
\]
since we must choose which \(k\) of the \(n\) positions are successes.  
Therefore, the PMF of \(X\) is
\[
\mathbb{P}(X = k) = \binom{n}{k} p^k (1-p)^{n-k}, \quad k = 0,1,\dots,n.
\]

\bigskip\textsc{Interpretation and shape}

When the coin is fair, \(p = 1/2\), the PMF is largest near \(k = n/2\).  
For example, with \(n = 10\), the most likely outcome is \(X = 5\), and probabilities decrease as \(k\) moves away from \(5\) in either direction.  
For larger \(n\), such as \(n = 100\), the most likely values of \(X\) form a band around \(np\); values far from \(np\) have very small probability.

When \(p\) is small (for instance \(p = 0.2\)), small values of \(X\) are more likely.  
If \(n\) is moderate, such as \(n = 10\), we typically see a few successes, often between \(0\) and \(4\).  
If \(n\) is large, such as \(n = 100\) with \(p = 0.1\), we say we expect “about” \(np = 10\) successes, meaning that values near \(10\) (roughly in some neighborhood such as from \(0\) up to around \(20\)) have noticeable probability, while much larger values are extremely unlikely.


\subsubsection{Geometric random variables}

A \textit{geometric} random variable models the number of trials needed until the first success occurs, in a sequence of independent, identical trials with success probability \(p\) at each trial.  
Typical examples are repeated coin tosses until the first head appears, or repeated attempts at an experiment until it succeeds.

\bigskip\textsc{Definition and PMF}

Consider an infinite sequence of independent trials, each with probability \(p\) of success and \(1-p\) of failure, where \(0 < p \le 1\).  
Let \(X\) be the number of the trial on which the first success occurs.  
Then \(X\) is called a \textit{geometric} random variable.  
Its possible values are the positive integers \(1,2,3,\dots\).

To compute the PMF, fix \(k \ge 1\).  
The event \(\{X = k\}\) means that the first \(k-1\) trials are failures and the \(k\)-th trial is a success.  
By independence,
\[
\mathbb{P}(X = k) = (1-p)^{k-1} p, \qquad k = 1,2,3,\dots
\]
As \(k\) increases, the factor \((1-p)^{k-1}\) makes these probabilities decrease geometrically.

\bigskip\textsc{Remark on “never seeing a success”}

There is a logically possible outcome where every trial is a failure (for example, all tails when tossing a coin), so there is no first success and \(X\) would be “infinite”.  
However, the probability of this event is
\[
\lim_{k \to \infty} (1-p)^k = 0,
\]
because \(0 < 1-p < 1\).  
So with probability \(1\) the random variable \(X\) takes a finite value in \(\{1,2,3,\dots\}\), and the probabilities \((1-p)^{k-1}p\) sum to \(1\).



\subsubsection{Expectation}

The \textit{expected value} (or expectation, or mean) of a random variable is a single number that describes what value we see “on average” in many independent repetitions of the experiment.  
For a discrete random variable, it is defined as a weighted average of its possible values, where the weights are the corresponding probabilities.

\bigskip\textsc{Definition for a discrete random variable}

Let \(X\) be a discrete random variable with PMF \(p_X(x) = \mathbb{P}(X = x)\).  
The expected value of \(X\), denoted \(\mathbb{E}[X]\), is
\[
\mathbb{E}[X] = \sum_{x} x \, p_X(x),
\]
where the sum runs over all possible values of \(X\).  
Each term is “value \(\times\) probability of that value”.

If \(X\) takes infinitely many (discrete) values, this becomes an infinite series.  
We then assume that
\[
\sum_{x} |x| \, p_X(x) < \infty,
\]
so that \(\mathbb{E}[X]\) is well defined and finite.

\bigskip\textsc{Example: Game of chance}

Suppose a random variable \(X\) represents your gain in a single play of a game, with
\[
\mathbb{P}(X = 1) = \frac{2}{10}, \quad
\mathbb{P}(X = 2) = \frac{5}{10}, \quad
\mathbb{P}(X = 4) = \frac{3}{10}.
\]
Then
\[
\mathbb{E}[X] = 1 \cdot \frac{2}{10}
             + 2 \cdot \frac{5}{10}
             + 4 \cdot \frac{3}{10}
           = \frac{2}{10} + \frac{10}{10} + \frac{12}{10}
           = \frac{24}{10}
           = 2.4.
\]
If you play many times, your average gain per game will be close to \(2.4\).

\bigskip\textsc{Bernoulli and indicator random variables}

Let \(X\) be a Bernoulli random variable with parameter \(p\), so
\[
\mathbb{P}(X = 1) = p, \qquad \mathbb{P}(X = 0) = 1 - p.
\]
Then
\[
\mathbb{E}[X] = 1 \cdot p + 0 \cdot (1-p) = p.
\]

As a special case, let \(X = I_A\) be the indicator of an event \(A\), so that \(X = 1\) if \(A\) occurs and \(X = 0\) otherwise.  
Then
\[
\mathbb{P}(I_A = 1) = \mathbb{P}(A),
\]
and
\[
\mathbb{E}[I_A] = \mathbb{P}(A).
\]
So the expected value of an indicator random variable equals the probability of the corresponding event.

\bigskip\textsc{Uniform random variable on \(\{0,1,\dots,n\}\)}

Let \(X\) be discrete uniform on \(\{0,1,\dots,n\}\).  
Then
\[
\mathbb{P}(X = k) = \frac{1}{n+1}, \quad k = 0,1,\dots,n.
\]
Using the definition,
\[
\mathbb{E}[X] = \sum_{k=0}^{n} k \cdot \frac{1}{n+1}
              = \frac{1}{n+1} \sum_{k=0}^{n} k
              = \frac{1}{n+1} \cdot \frac{n(n+1)}{2}
              = \frac{n}{2}.
\]
Thus the expectation is the midpoint of the interval \(\{0,\dots,n\}\).  
More generally, when a PMF is symmetric, the expected value lies at the center of symmetry.

\bigskip\textsc{Population-average interpretation}

Consider a class of \(n\) students, and let \(x_i\) be the weight of the \(i\)-th student.  
You perform an experiment where you pick one student uniformly at random, and let \(X\) be the weight of the selected student.  

The PMF of \(X\) is
\[
\mathbb{P}(X = x_i) = \frac{1}{n}, \quad i = 1,\dots,n.
\]
Then
\[
\mathbb{E}[X] = \sum_{i=1}^{n} x_i \cdot \frac{1}{n}
              = \frac{1}{n} \sum_{i=1}^{n} x_i.
\]
So in this example, the expected value of \(X\) is exactly the usual average (arithmetic mean) of the students’ weights.  
This illustrates that expectation can be seen both as the long-run average over many repetitions, and as an average over a finite population when all individuals are equally likely to be chosen.


\bigskip\textsc{Elementary properties of expectation}

A random variable \(X\) can be summarized by its expected value \(\mathbb{E}[X]\).  
Some basic properties of expectation are very natural and are worth stating explicitly.

\bigskip\textsc{Nonnegativity}

If \(X(\omega) \ge 0\) for every outcome \(\omega\) in the sample space, then
\[
\mathbb{E}[X] \ge 0.
\]
In words, if the random variable never takes negative values, its average value cannot be negative.  
This follows because \(\mathbb{E}[X]\) is a sum of terms of the form \(x \, p_X(x)\), and both factors are nonnegative.

\bigskip\textsc{Bounds from the range}

Suppose \(X(\omega)\) always lies between two constants \(a\) and \(b\), that is
\[
a \le X(\omega) \le b \quad \text{for all } \omega.
\]
Then the expectation also lies between \(a\) and \(b\):
\[
a \le \mathbb{E}[X] \le b.
\]

To see why \(\mathbb{E}[X] \ge a\), write
\[
\mathbb{E}[X] = \sum_{x} x \, p_X(x).
\]
Each possible value \(x\) of \(X\) satisfies \(x \ge a\), so
\[
\mathbb{E}[X]
= \sum_{x} x \, p_X(x)
\ge \sum_{x} a \, p_X(x)
= a \sum_{x} p_X(x)
= a,
\]
because the probabilities \(p_X(x)\) sum to \(1\).  
A symmetric argument with \(x \le b\) gives \(\mathbb{E}[X] \le b\).

\bigskip\textsc{Expectation of a constant}

A constant \(c\) can be viewed as a (degenerate) random variable that always takes the value \(c\).  
Its PMF assigns probability \(1\) to the value \(c\) and probability \(0\) to all other values.  
Then
\[
\mathbb{E}[c] = c \cdot 1 = c.
\]
So if a quantity is always equal to \(c\), its average is also \(c\).


\bigskip\textsc{Expected value rule}

A random variable \(Y\) is often defined as a function of another random variable \(X\), say
\[
Y = g(X).
\]
We want a convenient way to compute \(\mathbb{E}[Y]\) without first finding the PMF of \(Y\).  
The \textit{expected value rule} gives exactly this.


Let \(X\) be a discrete random variable with PMF \(p_X(x) = \mathbb{P}(X = x)\), and let
\[
Y = g(X)
\]
for some function \(g\).  
Then
\[
\mathbb{E}[Y] = \mathbb{E}[g(X)] = \sum_{x} g(x)\,p_X(x),
\]
where the sum is over all possible values of \(X\).  

Conceptually, we sum over the values of \(X\), and for each possible \(x\) we take the corresponding value \(g(x)\) of \(Y\), weighted by the probability that \(X\) equals \(x\).

\bigskip\textsc{Example}

Suppose \(X\) takes values \(2,3,4,5\) with probabilities
\[
\mathbb{P}(X=2)=0.1, \quad
\mathbb{P}(X=3)=0.2, \quad
\mathbb{P}(X=4)=0.3, \quad
\mathbb{P}(X=5)=0.4.
\]
Define \(Y = g(X)\) by
\[
g(2) = 3, \quad g(3) = 3, \quad g(4) = 4, \quad g(5) = 4.
\]
Then
\[
\mathbb{E}[Y]
= \sum_{x} g(x)\,p_X(x)
= 3 \cdot 0.1 + 3 \cdot 0.2 + 4 \cdot 0.3 + 4 \cdot 0.4.
\]
This matches the intuitive reasoning:  
10% of the time \(X=2\) and \(Y=3\), 20% of the time \(X=3\) and \(Y=3\), 30% of the time \(X=4\) and \(Y=4\), and 40% of the time \(X=5\) and \(Y=4\).

\bigskip\textsc{Why the rule works (idea)}

By definition,
\[
\mathbb{E}[Y] = \sum_{y} y\,\mathbb{P}(Y=y).
\]
For each fixed \(y\), the event \(\{Y=y\}\) is the union of all outcomes for which \(X\) takes some value \(x\) with \(g(x)=y\).  
Thus
\[
\mathbb{P}(Y=y) = \sum_{x : g(x)=y} \mathbb{P}(X=x).
\]
Substituting this into the definition and rearranging the order of summation leads to
\[
\mathbb{E}[Y]
= \sum_{y} y \sum_{x : g(x)=y} \mathbb{P}(X=x)
= \sum_{x} g(x)\,\mathbb{P}(X=x),
\]
which is the expected value rule.

\bigskip\textsc{A useful special case}

If \(g(x) = x^2\), then
\[
\mathbb{E}[X^2] = \sum_{x} x^2\,p_X(x).
\]
In general,
\[
\mathbb{E}[g(X)] \ne g(\mathbb{E}[X])
\]
for most functions \(g\).  
So, for example, typically
\[
\mathbb{E}[X^2] \ne (\mathbb{E}[X])^2.
\]


\bigskip\textsc{Linearity for one random variable}

Linearity of expectation describes how expectation behaves under multiplication by constants and addition.  
It says that taking expectations is a linear operation.


Let \(X\) be a random variable and let \(a, b\) be constants.  
Define a new random variable
\[
Y = aX + b.
\]
Then
\[
\mathbb{E}[aX + b] = a\,\mathbb{E}[X] + b.
\]

\bigskip\textsc{Derivation}

We view \(Y\) as a function of \(X\), namely \(Y = g(X)\) with \(g(x) = ax + b\).  
By the expected value rule,
\[
\mathbb{E}[Y] = \mathbb{E}[g(X)] = \sum_{x} g(x)\,p_X(x)
              = \sum_{x} (ax + b)\,p_X(x).
\]
We split the sum:
\[
\mathbb{E}[Y]
= \sum_{x} ax\,p_X(x) + \sum_{x} b\,p_X(x)
= a \sum_{x} x\,p_X(x) + b \sum_{x} p_X(x).
\]
The first sum is \(\mathbb{E}[X]\).  
The second sum is \(1\), since the probabilities \(p_X(x)\) add up to \(1\).  
Thus
\[
\mathbb{E}[aX + b] = a\,\mathbb{E}[X] + b.
\]

\bigskip\textsc{Special case: average salary}

If \(X\) is the salary of a random person in a population, then \(\mathbb{E}[X]\) is the average salary.  
If everyone’s salary is updated to
\[
Y = 2X + 100,
\]
then the new average is
\[
\mathbb{E}[Y] = 2\,\mathbb{E}[X] + 100.
\]
Doubling all salaries and adding \(100\) to each one doubles the average and adds \(100\).

\bigskip\textsc{Comment}

For linear functions \(g(x) = ax + b\), we have
\[
\mathbb{E}[g(X)] = g(\mathbb{E}[X]).
\]
For general, nonlinear functions \(g\), this identity usually does \textit{not} hold.



\subsubsection{Variance}

Variance measures how much a random variable tends to deviate (in squared distance) from its mean.  
It complements the expectation by quantifying the spread of the distribution around the average.

\bigskip\textsc{Definition}

Let \(X\) be a random variable with mean \(\mu = \mathbb{E}[X]\).  
The \textit{variance} of \(X\) is defined as
\[
\operatorname{Var}(X) = \mathbb{E}\bigl[(X - \mu)^2\bigr].
\]
This is the expected value of the squared distance from the mean.  
Since \((X - \mu)^2 \ge 0\) always, the variance is always nonnegative.

Using the expected value rule, for a discrete random variable with PMF \(p_X(x)\), we can write
\[
\operatorname{Var}(X)
= \sum_{x} (x - \mu)^2\,p_X(x),
\]
where the sum is over all possible values of \(X\). 
For each value \(x\), we compute its squared distance from \(\mu\) and weight it by the probability that \(X\) takes that value.

\bigskip\textsc{Standard deviation}

The \textit{standard deviation} of \(X\) is defined as
\[
\sigma_X = \sqrt{\operatorname{Var}(X)}.
\]
It has the same units as \(X\) (unlike the variance, which has squared units), and is often easier to interpret as a measure of typical deviation from the mean. 

\bigskip\textsc{Effect of adding and scaling}

Let \(Y = aX + b\), where \(a\) and \(b\) are constants.  
Then
\[
\operatorname{Var}(Y) = a^2 \operatorname{Var}(X).
\]
In particular,
\[
\operatorname{Var}(X + b) = \operatorname{Var}(X), \qquad
\operatorname{Var}(aX) = a^2 \operatorname{Var}(X) [][].
\]

Intuitively, adding a constant \(b\) just shifts the distribution right or left without changing its spread, so the variance stays the same.  
Multiplying by \(a\) stretches or shrinks the distribution horizontally by a factor \(|a|\), so the spread (variance) scales by \(a^2\)  

\bigskip\textsc{Example}

If \(\operatorname{Var}(X) = 2\), then
\[
\operatorname{Var}(3 - 4X) = (-4)^2 \operatorname{Var}(X) = 16 \cdot 2 = 32.
\]
The additive constant \(3\) does not affect the variance, while the factor \(-4\) multiplies the variance by \(16\).

\bigskip\textsc{Alternative formula}

There is a useful identity that often simplifies calculations:
\[
\operatorname{Var}(X) = \mathbb{E}[X^2] - \bigl(\mathbb{E}[X]\bigr)^2 [][].
\]
It comes from expanding
\[
\operatorname{Var}(X)
= \mathbb{E}\bigl[(X - \mu)^2\bigr]
= \mathbb{E}\bigl[X^2 - 2\mu X + \mu^2\bigr]
= \mathbb{E}[X^2] - 2\mu \mathbb{E}[X] + \mu^2
= \mathbb{E}[X^2] - \mu^2,
\]
since \(\mu = \mathbb{E}[X]\).  
This formula allows us to compute the variance by first finding \(\mathbb{E}[X]\) and \(\mathbb{E}[X^2]\), and then taking their difference.




\bigskip\textsc{Variance of a Bernoulli random variable}

A \textit{Bernoulli} random variable and a \textit{discrete uniform} random variable are simple examples where the variance can be computed explicitly.  
We first recall the Bernoulli case, then the uniform on \(\{0,\dots,n\}\), and finally a general uniform on \(\{a,\dots,b\}\). 

Let \(X\) be Bernoulli with parameter \(p\), so
\[
\mathbb{P}(X=1) = p, \qquad \mathbb{P}(X=0) = 1-p.
\]
We know that
\[
\mathbb{E}[X] = p.
\]
Using the alternative variance formula,
\[
\operatorname{Var}(X) = \mathbb{E}[X^2] - \bigl(\mathbb{E}[X]\bigr)^2 [][].
\]
For a Bernoulli random variable, \(X^2 = X\), since \(0^2=0\) and \(1^2=1\).  
Hence
\[
\mathbb{E}[X^2] = \mathbb{E}[X] = p,
\]
and
\[
\operatorname{Var}(X) = p - p^2 = p(1-p).
\]

As a function of \(p\), \(p(1-p)\) is a parabola on \([0,1]\) with maximum at \(p=\tfrac12\) and zeros at \(p=0\) and \(p=1\). 
This matches the intuition that a fair coin (\(p=\tfrac12\)) is “most random”, while a deterministic coin (\(p=0\) or \(p=1\)) has variance \(0\).

\bigskip\textsc{Variance of a uniform on \(\{0,1,\dots,n\}\)}

Let \(X\) be discrete uniform on \(\{0,1,\dots,n\}\).  
Then
\[
\mathbb{P}(X=k) = \frac{1}{n+1}, \quad k=0,1,\dots,n.
\]

We first compute \(\mathbb{E}[X]\) and \(\mathbb{E}[X^2]\).

For the mean,
\[
\mathbb{E}[X]
= \sum_{k=0}^{n} k \cdot \frac{1}{n+1}
= \frac{1}{n+1} \sum_{k=0}^{n} k
= \frac{1}{n+1} \cdot \frac{n(n+1)}{2}
= \frac{n}{2} [][].
\]

For \(\mathbb{E}[X^2]\),
\[
\mathbb{E}[X^2]
= \sum_{k=0}^{n} k^2 \cdot \frac{1}{n+1}
= \frac{1}{n+1} \sum_{k=0}^{n} k^2.
\]
The sum of squares satisfies
\[
\sum_{k=0}^{n} k^2 = \frac{n(n+1)(2n+1)}{6} [][],
\]
so
\[
\mathbb{E}[X^2]
= \frac{1}{n+1} \cdot \frac{n(n+1)(2n+1)}{6}
= \frac{n(2n+1)}{6}.
\]

Now apply \(\operatorname{Var}(X) = \mathbb{E}[X^2] - \bigl(\mathbb{E}[X]\bigr)^2\):
\[
\operatorname{Var}(X)
= \frac{n(2n+1)}{6} - \left(\frac{n}{2}\right)^2
= \frac{n(2n+1)}{6} - \frac{n^2}{4}.
\]
Bringing to a common denominator,
\[
\operatorname{Var}(X)
= \frac{4n(2n+1) - 3n^2}{12}
= \frac{8n^2 + 4n - 3n^2}{12}
= \frac{5n^2 + 4n}{12}.
\]
This is the variance of the uniform distribution on \(\{0,1,\dots,n\}\) written via the sum-of-squares formula.

In many references, one instead writes the variance of the uniform on \(\{0,1,\dots,n\}\) as
\[
\operatorname{Var}(X) = \frac{n(n+2)}{12},
\]
which is algebraically equivalent after simplifying the sum carefully. 

\bigskip\textsc{Uniform on \(\{a,a+1,\dots,b\}\)}

Now consider a general discrete uniform random variable \(Y\) on \(\{a,a+1,\dots,b\}\).  
Let
\[
n = b-a.
\]
We can write
\[
Y = a + X,
\]
where \(X\) is uniform on \(\{0,1,\dots,n\}\).  
Adding the constant \(a\) does not change the variance, so
\[
\operatorname{Var}(Y) = \operatorname{Var}(X) [][].
\]
Thus, substituting \(n = b-a\),
\[
\operatorname{Var}(Y)
= \frac{(b-a)\bigl((b-a)+2\bigr)}{12},
\]
which expresses the variance in terms of the endpoints of the range. 


\newpage
\bigskip\textsc{Conditional PMF given an event}

Let \(X\) be a discrete random variable with PMF \(p_X(x) = \mathbb{P}(X = x)\).  
Let \(A\) be an event with \(\mathbb{P}(A) > 0\).  
The \textit{conditional PMF} of \(X\) given \(A\) is
\[
p_{X\mid A}(x) = \mathbb{P}(X = x \mid A)
= \frac{\mathbb{P}(X = x \cap A)}{\mathbb{P}(A)}.
\]
This is an ordinary PMF on the same set of values as \(X\), but computed in the conditional model where \(A\) is known to have occurred.  
Its entries satisfy
\[
\sum_{x} p_{X\mid A}(x) = 1.
\]

\bigskip\textsc{Conditional expectation given an event}

The \textit{conditional expectation} of \(X\) given \(A\) is the mean of \(X\) in the conditional model.  
It is defined by
\[
\mathbb{E}[X \mid A] = \sum_{x} x\,\mathbb{P}(X = x \mid A)
= \sum_{x} x\,p_{X\mid A}(x).
\]
This is the same formula as for an ordinary expectation, but with conditional probabilities throughout.

\bigskip\textsc{Example: Conditioning a uniform}

Let \(X\) be uniform on \(\{1,2,3,4\}\).  
Then
\[
\mathbb{P}(X=x) = \frac{1}{4}, \quad x=1,2,3,4.
\]
By symmetry,
\[
\mathbb{E}[X] = \frac{1+4}{2} = \frac{5}{2},
\]
and using the uniform-variance formula on \(\{1,2,3,4\}\),
\[
\operatorname{Var}(X) = \frac{1}{12}\,(4-1)\bigl((4-1)+2\bigr)
= \frac{1}{12} \cdot 3 \cdot 5 = \frac{5}{4}.
\]

Now condition on the event
\[
A = \{X \in \{2,3,4\}\}.
\]
In the conditional model, outcome \(1\) is impossible, so
\[
\mathbb{P}(X=1 \mid A) = 0.
\]
The remaining values \(2,3,4\) had equal probabilities in the original model, so they remain equally likely given \(A\).  
Since their probabilities must sum to \(1\),
\[
\mathbb{P}(X=2 \mid A)
= \mathbb{P}(X=3 \mid A)
= \mathbb{P}(X=4 \mid A)
= \frac{1}{3}.
\]
Thus the conditional PMF is uniform on \(\{2,3,4\}\).

The conditional expectation is the mean of this conditional uniform:
\[
\mathbb{E}[X \mid A] = \frac{2+4}{2} = 3.
\]

For the conditional variance \(\operatorname{Var}(X \mid A)\), we can compute directly.  
The mean is \(3\), so
\[
\operatorname{Var}(X \mid A)
= \mathbb{E}\bigl[(X-3)^2 \mid A\bigr]
= (2-3)^2 \cdot \frac{1}{3}
  + (3-3)^2 \cdot \frac{1}{3}
  + (4-3)^2 \cdot \frac{1}{3}
= 1 \cdot \frac{1}{3} + 0 \cdot \frac{1}{3} + 1 \cdot \frac{1}{3}
= \frac{2}{3}.
\]

Notice that
\[
\operatorname{Var}(X \mid A) = \frac{2}{3} < \operatorname{Var}(X) = \frac{5}{4}.
\]
In the conditional model we have more information (we know that \(X\) is not \(1\)), so there is less uncertainty, and the variance decreases.



\subsubsection{Total expectation theorem}

The Total Expectation Theorem provides a way to compute an expectation by splitting the probability space into scenarios and averaging the conditional expectations over those scenarios.  
It is the expectation analogue of the Total Probability Theorem.

Let \(X\) be a random variable and let \(A_1, A_2, A_3\) be events that form a partition of the sample space, that is:
\[
A_i \cap A_j = \emptyset \ \text{for } i \ne j, 
\quad A_1 \cup A_2 \cup A_3 = \Omega,
\quad \mathbb{P}(A_i) > 0.
\]
Then the \textit{Total Expectation Theorem} says that
\[
\mathbb{E}[X]
= \mathbb{P}(A_1)\,\mathbb{E}[X \mid A_1]
+ \mathbb{P}(A_2)\,\mathbb{E}[X \mid A_2]
+ \mathbb{P}(A_3)\,\mathbb{E}[X \mid A_3].
\]
More generally, for a finite or countable partition \(\{A_i\}\),
\[
\mathbb{E}[X] = \sum_i \mathbb{P}(A_i)\,\mathbb{E}[X \mid A_i].
\]

\bigskip\textsc{From total probability to total expectation}

For a fixed value \(x\), the Total Probability Theorem applied to the event \(\{X = x\}\) gives
\[
\mathbb{P}(X = x)
= \sum_i \mathbb{P}(X = x \mid A_i)\,\mathbb{P}(A_i).
\]
This holds for every possible value \(x\).  
Multiply both sides by \(x\) and sum over all \(x\):
\[
\sum_x x\,\mathbb{P}(X = x)
= \sum_x x \sum_i \mathbb{P}(X = x \mid A_i)\,\mathbb{P}(A_i).
\]
The left-hand side is \(\mathbb{E}[X]\).  
On the right-hand side, we can interchange the order of summation:
\[
\mathbb{E}[X]
= \sum_i \mathbb{P}(A_i) \sum_x x\,\mathbb{P}(X = x \mid A_i).
\]
But for each \(i\),
\[
\sum_x x\,\mathbb{P}(X = x \mid A_i)
= \mathbb{E}[X \mid A_i].
\]
Substituting this into the previous expression gives the Total Expectation Theorem.

\bigskip\textsc{Interpretation}

We can think of \(\{A_1, A_2, A_3\}\) as different scenarios.  
Under each scenario \(A_i\), the random variable \(X\) has a conditional distribution and a conditional expectation \(\mathbb{E}[X \mid A_i]\).  
The theorem says that the overall expectation \(\mathbb{E}[X]\) is the weighted average of these conditional expectations, where the weights are the scenario probabilities \(\mathbb{P}(A_i)\).

\bigskip\textsc{Example: Splitting a PMF into two scenarios}

Consider a discrete random variable \(X\) with PMF having six mass points: three “small” values and three “large” values.  
Define
\[
A_1 = \{\text{“small” values of } X\}, \quad
A_2 = \{\text{“large” values of } X\},
\]
so that \(\{A_1,A_2\}\) is a partition of the sample space.

Suppose the probabilities of the three small values are all \(1/9\), so
\[
\mathbb{P}(A_1) = \frac{1}{9} + \frac{1}{9} + \frac{1}{9} = \frac{1}{3},
\]
and the probabilities of the three large values are all \(2/9\), so
\[
\mathbb{P}(A_2) = \frac{2}{9} + \frac{2}{9} + \frac{2}{9} = \frac{2}{3}.
\]

In the conditional model given \(A_1\), the three small values are equally likely, so \(X\) is uniform on \(\{0,1,2\}\) and
\[
\mathbb{E}[X \mid A_1] = 1
\]
(the midpoint).  
Similarly, given \(A_2\), \(X\) is uniform on \(\{6,7,8\}\), so
\[
\mathbb{E}[X \mid A_2] = 7.
\]

By the Total Expectation Theorem,
\[
\mathbb{E}[X]
= \mathbb{P}(A_1)\,\mathbb{E}[X \mid A_1]
+ \mathbb{P}(A_2)\,\mathbb{E}[X \mid A_2]
= \frac{1}{3} \cdot 1 + \frac{2}{3} \cdot 7
= \frac{1}{3} + \frac{14}{3}
= 5.
\]
Thus, instead of computing \(\mathbb{E}[X]\) directly from all six probabilities at once, we divided the problem into two simpler conditional expectations and combined them with appropriate weights.



\subsubsection{Geometric PMF and memorylessness}

Let \(X\) be the number of tosses until the first heads.  
Then \(X\) takes values \(1,2,3,\dots\), and
\[
\mathbb{P}(X = k) = (1-p)^{k-1} p, \quad k = 1,2,\dots
\]

The key property is \textit{memorylessness}.  
Suppose we know that the first toss was tails, so \(X > 1\).  
From that point on, we are again waiting for the first heads in a sequence of independent tosses with success probability \(p\).  
Thus, the number of \textit{additional} tosses until heads, namely \(X-1\), has the same geometric distribution as \(X\), when we condition on \(X>1\):
\[
X-1 \mid (X>1) \sim \text{Geometric}(p).
\]
More generally, for any integer \(n \ge 1\),
\[
X-n \mid (X>n) \sim \text{Geometric}(p).
\]

Concretely, for example
\[
\mathbb{P}(X-1 = 3 \mid X>1)
= \mathbb{P(\text{T,T,T,H on tosses }2,3,4 \mid \text{T on toss }1)}.
\]
Because tosses are independent, conditioning on tails on the first toss does not change the probabilities of future sequences, so this equals
\[
\mathbb{P}(\text{T,T,T,H on tosses }2,3,4)
= (1-p)^3 p,
\]
which is exactly \(\mathbb{P}(X = 4)\), the geometric PMF at \(k=4\).  
This illustrates that the conditional distribution of the remaining waiting time is geometric with parameter \(p\).

\bigskip\textsc{Using memorylessness to find the expectation}

Directly computing
\[
\mathbb{E}[X] = \sum_{k=1}^{\infty} k (1-p)^{k-1} p
\]
is possible but involves an infinite series.  
Instead, we use memorylessness and the total expectation idea.

Write
\[
X = 1 + (X-1).
\]
Taking expectations,
\[
\mathbb{E}[X] = 1 + \mathbb{E}[X-1].
\]

Now split \(\mathbb{E}[X-1]\) into two scenarios according to the outcome of the first toss:
\[
\mathbb{E}[X-1]
= \mathbb{E}[X-1 \mid X=1]\,\mathbb{P}(X=1)
+ \mathbb{E}[X-1 \mid X>1]\,\mathbb{P}(X>1).
\]
If \(X=1\), then \(X-1=0\), so
\[
\mathbb{E}[X-1 \mid X=1] = 0, \quad \mathbb{P}(X=1) = p.
\]
Thus the first term is zero.  
The second term is
\[
\mathbb{E}[X-1 \mid X>1]\,\mathbb{P}(X>1)
= \mathbb{E}[X-1 \mid X>1] (1-p).
\]

By memorylessness,
\[
X-1 \mid (X>1) \sim \text{Geometric}(p),
\]
so
\[
\mathbb{E}[X-1 \mid X>1] = \mathbb{E}[X].
\]
Therefore
\[
\mathbb{E}[X-1] = (1-p)\,\mathbb{E}[X].
\]

Substitute back into \(\mathbb{E}[X] = 1 + \mathbb{E}[X-1]\):
\[
\mathbb{E}[X] = 1 + (1-p)\,\mathbb{E}[X].
\]
Rearrange:
\[
\mathbb{E}[X] - (1-p)\,\mathbb{E}[X] = 1,
\]
\[
p\,\mathbb{E}[X] = 1,
\]
\[
\mathbb{E}[X] = \frac{1}{p}.
\]

This makes intuitive sense: if the success probability \(p\) is small, we expect to wait longer on average until the first heads.



\subsubsection{Joint PMFs and the expected value rule}

When dealing with several discrete random variables together, their joint behaviour is described by a \textit{joint} PMF.  
From the joint PMF we can recover the individual (marginal) PMFs and compute expectations of functions of several variables.

\bigskip\textsc{Joint and marginal PMFs}

Let \(X\) and \(Y\) be discrete random variables.  
Their \textit{joint} PMF is
\[
p_{X,Y}(x,y) = \mathbb{P}(X = x,\, Y = y).
\]
This is a function of two arguments: for each pair \((x,y)\) it gives the probability that \(X\) takes value \(x\) and \(Y\) takes value \(y\) simultaneously.  
It satisfies
\[
\sum_{x}\sum_{y} p_{X,Y}(x,y) = 1.
\]

The \textit{marginal} PMFs of \(X\) and \(Y\) are obtained by summing out the other variable:
\[
p_X(x) = \mathbb{P}(X = x) = \sum_{y} p_{X,Y}(x,y),
\]
\[
p_Y(y) = \mathbb{P}(Y = y) = \sum_{x} p_{X,Y}(x,y).
\]
For example, if the joint PMF is given in a table, then \(p_X(x)\) is the sum of the entries in the column corresponding to \(x\), and \(p_Y(y)\) is the sum of the entries in the row corresponding to \(y\).

\bigskip\textsc{Example: reading a joint PMF table}

Suppose the joint PMF of \(X\) and \(Y\), for \(x \in \{0,1,2\}\) and \(y \in \{1,2,3\}\), is given in a \(3\times 3\) table of probabilities.  
Then:
- The probability that \(X=1\) and \(Y=3\) is the entry in the row \(y=3\), column \(x=1\).  
- The marginal \(\mathbb{P}(X=2)\) is the sum of the entries in the column \(x=2\).  
- The marginal \(\mathbb{P}(Y=2)\) is the sum of the entries in the row \(y=2\).  

If we want \(\mathbb{P}(X=Y)\), we add the probabilities of all diagonal pairs \((0,0)\), \((1,1)\), \((2,2)\) (or whichever pairs satisfy \(x=y\)) in the joint table.

\bigskip\textsc{Joint PMF for more variables}

For three discrete random variables \(X,Y,Z\), the joint PMF is
\[
p_{X,Y,Z}(x,y,z) = \mathbb{P}(X=x, Y=y, Z=z),
\]
and
\[
\sum_{x}\sum_{y}\sum_{z} p_{X,Y,Z}(x,y,z) = 1.
\]
From this we can obtain marginals by summing out the variables we are not interested in, for example
\[
p_X(x) = \sum_{y}\sum_{z} p_{X,Y,Z}(x,y,z),
\]
\[
p_{X,Y}(x,y) = \sum_{z} p_{X,Y,Z}(x,y,z).
\]

\bigskip\textsc{Functions of two random variables and their PMF}

Given two random variables \(X\) and \(Y\), a function
\[
Z = g(X,Y)
\]
is itself a random variable.  
To find the PMF of \(Z\), we use the joint PMF of \((X,Y)\):
\[
p_Z(z) = \mathbb{P}(Z = z)
= \mathbb{P}\bigl(g(X,Y) = z\bigr)
= \sum_{(x,y):\, g(x,y) = z} p_{X,Y}(x,y).
\]
That is, we sum the probabilities of all pairs \((x,y)\) for which \(g(x,y)\) equals the value \(z\).

\bigskip\textsc{Expected value rule for functions of two variables}

For a function \(g(X,Y)\), there is an expected value rule analogous to the one-variable case:
\[
\mathbb{E}[g(X,Y)] = \sum_{x}\sum_{y} g(x,y)\,p_{X,Y}(x,y).
\]
Interpretation: for each possible pair \((x,y)\), the random variables take this pair with probability \(p_{X,Y}(x,y)\), and then the function has value \(g(x,y)\); the product \(g(x,y)\,p_{X,Y}(x,y)\) is the contribution of that pair to the overall average, and we sum over all pairs.


\subsubsection{Linearity of expectations}

Linearity of expectation extends to sums of multiple random variables and lets us compute the binomial mean very easily.

\bigskip\textsc{Linearity for several random variables}

For any discrete random variables \(X_1,\dots,X_n\),
\[
\mathbb{E}\Bigl[\sum_{i=1}^{n} X_i\Bigr] = \sum_{i=1}^{n} \mathbb{E}[X_i].
\]
For example, for three variables,
\[
\mathbb{E}[2X + 3Y - Z]
= 2\,\mathbb{E}[X] + 3\,\mathbb{E}[Y] - \mathbb{E}[Z].
\]

\bigskip\textsc{Derivation for two variables}

For two variables \(X\) and \(Y\),
\[
\mathbb{E}[X+Y]
= \sum_{x}\sum_{y} (x+y)\,p_{X,Y}(x,y).
\]
Split the sum:
\[
\mathbb{E}[X+Y]
= \sum_{x}\sum_{y} x\,p_{X,Y}(x,y)
+ \sum_{x}\sum_{y} y\,p_{X,Y}(x,y).
\]
In the first term, for fixed \(x\),
\[
\sum_{y} x\,p_{X,Y}(x,y)
= x \sum_{y} p_{X,Y}(x,y)
= x\,p_X(x),
\]
so
\[
\sum_{x}\sum_{y} x\,p_{X,Y}(x,y)
= \sum_{x} x\,p_X(x)
= \mathbb{E}[X].
\]
A symmetric argument shows that the second term equals \(\mathbb{E}[Y]\).  
Thus
\[
\mathbb{E}[X+Y] = \mathbb{E}[X] + \mathbb{E}[Y].
\]
The general \(n\)-variable case follows similarly.

\bigskip\textsc{Mean of a binomial via indicators}

Let \(X \sim \text{Binomial}(n,p)\).  
Interpret \(X\) as the number of successes in \(n\) independent trials, each with success probability \(p\).

Define indicator variables
\[
X_i =
\begin{cases}
1 & \text{if trial \(i\) is a success},\\
0 & \text{otherwise},
\end{cases}
\qquad i=1,\dots,n.
\]
Each \(X_i\) is Bernoulli\((p)\), so
\[
\mathbb{P}(X_i = 1) = p, \quad \mathbb{P}(X_i = 0) = 1-p, \quad \mathbb{E}[X_i] = p.
\]

The total number of successes is the sum of these indicators:
\[
X = X_1 + X_2 + \cdots + X_n.
\]
By linearity of expectation,
\[
\mathbb{E}[X]
= \mathbb{E}[X_1] + \cdots + \mathbb{E}[X_n]
= p + \cdots + p
= np.
\]

So a binomial random variable with parameters \(n\) and \(p\) has mean
\[
\mathbb{E}[X] = np,
\]
which matches the intuition that in \(n\) trials with success probability \(p\), we expect \(np\) successes on average.



\newpage
\subsubsection{Conditional PMFs}

When we condition on the value of another random variable, we obtain a family of conditional PMFs for \(X\), one for each value of \(Y\).

\bigskip\textsc{Conditional PMF given another variable}

Let \(X\) and \(Y\) be discrete random variables with joint PMF \(p_{X,Y}(x,y)\) and marginal \(p_Y(y) = \mathbb{P}(Y=y)\).  
For any \(y\) with \(p_Y(y) > 0\), the \textit{conditional PMF} of \(X\) given \(Y=y\) is
\[
p_{X\mid Y}(x\mid y)
= \mathbb{P}(X = x \mid Y = y)
= \frac{p_{X,Y}(x,y)}{p_Y(y)}.
\]
If we fix \(y\), this is a PMF in the variable \(x\), and its values satisfy
\[
\sum_x p_{X\mid Y}(x\mid y) = 1.
\]
Changing \(y\) changes the conditional PMF, so \(\{p_{X\mid Y}(\cdot\mid y)\}\) is a family of PMFs indexed by \(y\).

\bigskip\textsc{Example from a joint table}

Suppose the joint PMF of \((X,Y)\) is given by a table, and we condition on \(Y=2\).  
First compute
\[
p_Y(2) = \sum_x p_{X,Y}(x,2),
\]
by adding the entries in the row \(y=2\).  
Then, for each \(x\),
\[
p_{X\mid Y}(x\mid 2)
= \frac{p_{X,Y}(x,2)}{p_Y(2)}.
\]
For instance, if the row \(y=2\) has entries proportional to \(0,1,3,1\), their sum is \(5\), so the conditional PMF values are
\[
0,\ \frac{1}{5},\ \frac{3}{5},\ \frac{1}{5}.
\]
Geometrically, the conditional PMF for \(X\mid Y=2\) is the corresponding row of the joint PMF, rescaled so that its entries add to \(1\).

\bigskip\textsc{Multiplication rule in PMF form}

From the definition
\[
p_{X\mid Y}(x\mid y) = \frac{p_{X,Y}(x,y)}{p_Y(y)},
\]
we can rewrite
\[
p_{X,Y}(x,y)
= p_Y(y)\,p_{X\mid Y}(x\mid y),
\]
which is the multiplication rule in PMF notation.  
Symmetrically,
\[
p_{X,Y}(x,y)
= p_X(x)\,p_{Y\mid X}(y\mid x).
\]
Thus we can specify a joint distribution either directly via \(p_{X,Y}\) or indirectly by giving a marginal (say \(p_Y\)) and a conditional PMF (say \(p_{X\mid Y}\)).

\bigskip\textsc{More than two variables}

For three random variables \(X,Y,Z\), the conditional PMF of \(X\) given \(Y=y\) and \(Z=z\) is
\[
p_{X\mid Y,Z}(x\mid y,z)
= \mathbb{P}(X=x \mid Y=y, Z=z)
= \frac{p_{X,Y,Z}(x,y,z)}{p_{Y,Z}(y,z)},
\]
where
\[
p_{Y,Z}(y,z) = \sum_x p_{X,Y,Z}(x,y,z).
\]

We also have a multiplication rule for three variables.  
Let
\[
A = \{X=x\}, \quad B = \{Y=y\}, \quad C = \{Z=z\}.
\]
Then
\[
\mathbb{P}(A \cap B \cap C)
= \mathbb{P}(A)\,\mathbb{P}(B\mid A)\,\mathbb{P}(C\mid A \cap B),
\]
which in PMF notation becomes
\[
p_{X,Y,Z}(x,y,z)
= p_X(x)\,p_{Y\mid X}(y\mid x)\,p_{Z\mid X,Y}(z\mid x,y).
\]

\bigskip\textsc{conditional expectation}

A conditional expectation is just the expectation of \(X\) computed in the conditional model “given \(Y=y\)”.


Let \(X\) and \(Y\) be discrete random variables, and suppose \(\mathbb{P}(Y = y) > 0\).  
The conditional expectation of \(X\) given \(Y=y\) is
\[
\mathbb{E}[X \mid Y = y]
= \sum_{x} x \,\mathbb{P}(X = x \mid Y = y).
\]
Here we use the conditional PMF \(p_{X\mid Y}(x\mid y)\) instead of the original PMF of \(X\).

\bigskip\textsc{Expected value rule (conditional form)}

If \(g\) is a function and \(Z = g(X)\), then in the conditional model “given \(Y=y\)” we have
\[
\mathbb{E}[g(X) \mid Y = y]
= \sum_{x} g(x)\,\mathbb{P}(X = x \mid Y = y).
\]
This is the same expected value rule as before, but with conditional probabilities.

\bigskip\textsc{Total expectation theorem (version with \(Y\))}

Let \(X\) and \(Y\) be discrete random variables, and suppose \(Y\) takes values in some countable set \(\mathcal{Y}\).  
Then the total expectation theorem says
\[
\mathbb{E}[X]
= \sum_{y \in \mathcal{Y}} \mathbb{P}(Y = y)\,\mathbb{E}[X \mid Y = y].
\]
So we can view the possible values of \(Y\) as “scenarios”:  
for each scenario \(Y=y\), we compute the conditional mean \(\mathbb{E}[X \mid Y=y]\), and then average these means using \(\mathbb{P}(Y=y)\) as weights.

\bigskip\textsc{Remark on infinite ranges}

If \(Y\) can take infinitely many discrete values, the sum
\[
\sum_{y} \mathbb{P}(Y = y)\,\mathbb{E}[X \mid Y = y]
\]
becomes an infinite series.  
Under the usual assumption that \(\mathbb{E}[X]\) is well defined (finite), this identity still holds, and we can safely use the total expectation theorem even when \(Y\) ranges over an infinite discrete set.


\subsubsection{Independence of random variables}

Independence for random variables extends the idea of independence for events to the level of entire distributions.

Let \(X\) be a random variable and \(A\) an event.  
We say that \(X\) and \(A\) are \textit{independent} if, for every real number \(x\),
\[
\mathbb{P}(X = x \text{ and } A) = \mathbb{P}(X = x)\,\mathbb{P}(A).
\]
Equivalently, whenever \(\mathbb{P}(A) > 0\),
\[
\mathbb{P}(X = x \mid A) = \mathbb{P}(X = x) \quad \text{for all } x.
\]
Intuition: knowing that \(A\) occurred does not change the distribution of \(X\).  

Symmetrically, for any \(x\) with \(\mathbb{P}(X = x) > 0\),
\[
\mathbb{P}(A \mid X = x) = \mathbb{P}(A),
\]
so learning the value of \(X\) does not change your beliefs about \(A\).

\bigskip\textsc{Independence of two random variables}

Let \(X\) and \(Y\) be discrete random variables.  
We say that \(X\) and \(Y\) are \textit{independent} if, for all real \(x\) and \(y\),
\[
\mathbb{P}(X = x,\, Y = y)
= \mathbb{P}(X = x)\,\mathbb{P}(Y = y).
\]
In terms of PMFs,
\[
p_{X,Y}(x,y) = p_X(x)\,p_Y(y) \quad \text{for all } x,y.
\]
So the joint PMF factorizes as the product of the marginals.

Equivalently, for any \(y\) with \(\mathbb{P}(Y = y) > 0\),
\[
\mathbb{P}(X = x \mid Y = y) = \mathbb{P}(X = x) \quad \text{for all } x,
\]
that is,
\[
p_{X\mid Y}(x\mid y) = p_X(x).
\]
So knowing \(Y=y\) does not change the distribution of \(X\).  
Symmetrically,
\[
p_{Y\mid X}(y\mid x) = p_Y(y),
\]
so knowing \(X=x\) does not change the distribution of \(Y\).

\bigskip\textsc{Independence of multiple random variables}

For three random variables \(X,Y,Z\), we say they are \textit{mutually independent} if, for all \(x,y,z\),
\[
p_{X,Y,Z}(x,y,z)
= p_X(x)\,p_Y(y)\,p_Z(z).
\]
More generally, a finite collection \(X_1,\dots,X_n\) is mutually independent if the joint PMF factors as the product of all the marginals:
\[
p_{X_1,\dots,X_n}(x_1,\dots,x_n)
= \prod_{i=1}^{n} p_{X_i}(x_i)
\]
for all choices of \((x_1,\dots,x_n)\).

Intuitively, this means that information about some of the variables does not change the distribution of any of the others: all conditional PMFs coincide with the corresponding unconditional PMFs.  
In practical terms, independence models situations where each random variable is generated by its own separate probabilistic experiment, with no interaction and no shared source of randomness.



Independence gives one more important rule for expectations: products of independent random variables have factorizing expectations.

\bigskip\textsc{Product of independent random variables}

Let \(X\) and \(Y\) be independent discrete random variables.  
Then
\[
\mathbb{E}[XY] = \mathbb{E}[X]\,\mathbb{E}[Y].
\]

\textit{Derivation.}  
Using the expected value rule with the joint PMF:
\[
\mathbb{E}[XY]
= \sum_{x} \sum_{y} xy\,p_{X,Y}(x,y).
\]
Independence means
\[
p_{X,Y}(x,y) = p_X(x)\,p_Y(y),
\]
so
\[
\mathbb{E}[XY]
= \sum_{x} \sum_{y} xy\,p_X(x)\,p_Y(y).
\]
For fixed \(x\), the factor \(x\,p_X(x)\) does not depend on \(y\), hence
\[
\mathbb{E}[XY]
= \sum_{x} x\,p_X(x) \sum_{y} y\,p_Y(y)
= \Bigl(\sum_{x} x\,p_X(x)\Bigr)
  \Bigl(\sum_{y} y\,p_Y(y)\Bigr)
= \mathbb{E}[X]\,\mathbb{E}[Y].
\]

\bigskip\textsc{Products of functions of independent variables}

If \(X\) and \(Y\) are independent, and \(g\) and \(h\) are functions, then the random variables \(g(X)\) and \(h(Y)\) are also independent.  
Therefore,
\[
\mathbb{E}[g(X)\,h(Y)]
= \mathbb{E}[g(X)]\,\mathbb{E}[h(Y)].
\]
We can view this either as a consequence of independence of \(g(X)\) and \(h(Y)\), or prove it directly by repeating the previous derivation with \(g(x)\) in place of \(x\) and \(h(y)\) in place of \(y\).


We want to understand how variance behaves for sums, and then use this to find the variance of a binomial random variable.

\newpage
\bigskip\textsc{Variance of a sum under independence}

In general, \(\operatorname{Var}(X+Y)\) is not equal to \(\operatorname{Var}(X)+\operatorname{Var}(Y)\).  
However, if \(X\) and \(Y\) are independent, we do have
\[
\operatorname{Var}(X+Y) = \operatorname{Var}(X) + \operatorname{Var}(Y).
\]

Assume first that \(\mathbb{E}[X]=\mathbb{E}[Y]=0\), just to simplify the algebra.  
Then
\[
\operatorname{Var}(X+Y)
= \mathbb{E}\bigl[(X+Y)^2\bigr]
= \mathbb{E}\bigl[X^2 + 2XY + Y^2\bigr].
\]
Using linearity of expectation,
\[
\operatorname{Var}(X+Y)
= \mathbb{E}[X^2] + 2\mathbb{E}[XY] + \mathbb{E}[Y^2].
\]
Since the means are zero, \(\mathbb{E}[X^2] = \operatorname{Var}(X)\) and \(\mathbb{E}[Y^2] = \operatorname{Var}(Y)\).  

If \(X\) and \(Y\) are independent, then
\[
\mathbb{E}[XY] = \mathbb{E}[X]\,\mathbb{E}[Y] = 0 \cdot 0 = 0.
\]
So
\[
\operatorname{Var}(X+Y)
= \operatorname{Var}(X) + \operatorname{Var}(Y).
\]

Without the zero-mean assumption, one can expand \(\operatorname{Var}(X+Y)\) using the general formula \(\operatorname{Var}(Z)=\mathbb{E}[Z^2]-(\mathbb{E}[Z])^2\) and obtain the same result; the cross terms again cancel because \(\mathbb{E}[XY]=\mathbb{E}[X]\mathbb{E}[Y]\) under independence.

\bigskip\textsc{Examples showing dependence matters}

If \(Y = X\), then \(X+Y = 2X\), so
\[
\operatorname{Var}(X+Y) = \operatorname{Var}(2X) = 4\,\operatorname{Var}(X),
\]
which is generally not equal to \(\operatorname{Var}(X)+\operatorname{Var}(Y) = 2\,\operatorname{Var}(X)\).  
Here \(X\) and \(Y\) are completely dependent.

If \(Y = -X\), then \(X+Y = 0\) always, so
\[
\operatorname{Var}(X+Y) = \operatorname{Var}(0) = 0,
\]
again very different from \(\operatorname{Var}(X)+\operatorname{Var}(Y)\).  

These examples show that knowing just the individual variances is not enough; the dependence structure between \(X\) and \(Y\) is crucial.

\bigskip\textsc{Example with independence and scaling}

If \(X\) and \(Y\) are independent, then \(X\) is also independent of any function of \(Y\), such as \(-3Y\).  
Thus
\[
\operatorname{Var}(X - 3Y)
= \operatorname{Var}(X) + \operatorname{Var}(-3Y)
= \operatorname{Var}(X) + 9\,\operatorname{Var}(Y),
\]
using the rule \(\operatorname{Var}(aZ) = a^2 \operatorname{Var}(Z)\).

\bigskip\textsc{Variance of a binomial random variable}

Let \(X\) be binomial with parameters \(n\) and \(p\).  
Interpret \(X\) as the number of successes in \(n\) independent trials, each with success probability \(p\).

Define indicator random variables
\[
X_i =
\begin{cases}
1 & \text{if the \(i\)-th trial is a success},\\
0 & \text{otherwise},
\end{cases}
\quad i=1,\dots,n.
\]
Each \(X_i\) is Bernoulli\((p)\), and the total number of successes is
\[
X = X_1 + X_2 + \cdots + X_n.
\]
The independence of the trials means the indicators \(X_1,\dots,X_n\) are independent.  

For each \(X_i\),
\[
\mathbb{E}[X_i] = p, \qquad \operatorname{Var}(X_i) = p(1-p).
\]

Using independence and the variance-of-sum rule,
\[
\operatorname{Var}(X)
= \operatorname{Var}\bigl(X_1 + \cdots + X_n\bigr)
= \sum_{i=1}^{n} \operatorname{Var}(X_i)
= \sum_{i=1}^{n} p(1-p)
= n\,p(1-p).
\]

Thus a binomial random variable with parameters \(n\) and \(p\) has variance
\[
\operatorname{Var}(X) = n p (1-p).
\]

\subsubsection{The hat problem}


There are \(n\) people and \(n\) hats, all permutations equally likely.  
Let \(X\) be the number of people who get their own hat, and \(X_i\) the indicator for person \(i\).

\bigskip\textsc{Definition of the indicators}

For each \(i = 1,\dots,n\),
\[
X_i =
\begin{cases}
1 & \text{if person } i \text{ gets their own hat},\\
0 & \text{otherwise}.
\end{cases}
\]
Then
\[
X = X_1 + X_2 + \cdots + X_n.
\]

\bigskip\textsc{Expectation of each \(X_i\)}

In a random permutation, person \(i\) is equally likely to receive any of the \(n\) hats.  
Only one of these is their own hat, so
\[
\mathbb{P}(X_i = 1) = \frac{1}{n}.
\]
Thus
\[
\mathbb{E}[X_i]
= 1 \cdot \frac{1}{n} + 0 \cdot \Bigl(1 - \frac{1}{n}\Bigr)
= \frac{1}{n}.
\]

\bigskip\textsc{Expectation of \(X\)}

By linearity of expectation,
\[
\mathbb{E}[X]
= \sum_{i=1}^{n} \mathbb{E}[X_i]
= \sum_{i=1}^{n} \frac{1}{n}
= 1.
\]

So in the hat problem, the expected number of people who get their own hat back is always
\[
\mathbb{E}[X] = 1,
\]
independent of \(n\).